% Fichier complété par tools/complete_missing_corrections.py.
% Source : PPM/PPM-01/1005_Dessin2D/1005_Dessin2D.tex
% Statut : rédigé (1 question(s)).
\begin{corrigebox}[Corrigé rédigé]
\CorrigeQuestion{1}
On complète les vues par projection orthogonale : chaque sommet de la
        vue connue est projeté perpendiculairement vers la vue à compléter, les
        profondeurs étant reportées avec la ligne à 45 degrés. Les arêtes visibles sont en
        trait fort continu, les arêtes cachées en trait interrompu fin et les axes en trait
        mixte fin. On vérifie la correspondance stricte des largeurs entre face/dessus et
        des hauteurs entre face/profil. Les formes inclinées se projettent par leurs points
        extrêmes; les évidements sont reportés sur les trois vues. Le résultat final doit
        être cohérent avec la perspective fournie et ne contenir aucune arête redondante.
\end{corrigebox}
